% Part: set-theory
% Chapter: ordinals
% Section: opps

\documentclass[../../../include/open-logic-section]{subfiles}

\begin{document}

\olfileid{sth}{ordinals}{opps}
\olsection{উত্তরসূরি ও সীমা ক্রমসংখ্যা}

পরের কয়েকটি অধ্যায়ে আমরা ক্রমসংখ্যা বহুবার ব্যবহার
করব। তাই কয়েকটি সহজ ধারণা আগে প্রবর্তন করে নেওয়া
উপকারী।

\begin{defn}
যেকোনো ক্রমসংখ্যা $\alpha$-এর \emph{উত্তরসূরি} হলো $\ordsucc{\alpha}
=\alpha \cup \{\alpha\}$। যদি কোনো ক্রমসংখ্যা~$\beta$-র
জন্য $\ordsucc{\beta} = \alpha$ হয়, তবে $\alpha$-কে
একটি \emph{উত্তরসূরি} ক্রমসংখ্যা বলি। $\alpha$-কে
একটি \emph{সীমা} ক্রমসংখ্যা বলি যদি এবং কেবল যদি
$\alpha$ শূন্যও নয়, উত্তরসূরি ক্রমসংখ্যাও নয়।
\end{defn}
\noindent
নিচের ফলটি দেখায়, এটাই \emph{উত্তরসূরি}-র উপযুক্ত ধারণা:

\begin{prop}
যেকোনো ক্রমসংখ্যা $\alpha$-এর জন্য:
\begin{enumerate}
	\item $\alpha \in \ordsucc{\alpha}$;
	\item $\ordsucc{\alpha}$ একটি ক্রমসংখ্যা;
	\item $\alpha \in \beta \in
	\ordsucc{\alpha}$ এমন কোনো ক্রমসংখ্যা $\beta$ নেই।
	\end{enumerate}
\end{prop}

\begin{proof}
স্পষ্টতই $\alpha \in \alpha \cup \{\alpha\} = \ordsucc{\alpha}$।
একইভাবে, $\ordsucc{\alpha}$ ক্রমসংখ্যাগুলির একটি
পরিযায়ী সেট; তাই \olref[basic]{corordtransitiveord}
অনুযায়ী এটিও ক্রমসংখ্যা। আবার $\alpha \in \beta
\in \ordsucc{\alpha}$ অসম্ভব, কারণ তা হলে হয়
$\beta \in \alpha$, নয় $\beta = \alpha$;
উভয় ক্ষেত্রেই \olref[basic]{ordordered}-এর বিরোধ হয়।
\end{proof}
\noindent
এটি ট্রান্সফাইনাইট আবেশে প্রমাণের একটি ভিন্ন রূপও দেয়:

\begin{thm}[সরল ট্রান্সফাইনাইট আবেশ]\ollabel{simpletransrecursion}
ধরি, $\phi(x)$ এমন একটি সূত্র, যার জন্য:
\begin{enumerate}
	\item $\phi(\emptyset)$ সত্য; এবং
	\item যেকোনো ক্রমসংখ্যা $\alpha$-এর জন্য, $\phi(\alpha)$
	থেকে $\phi(\ordsucc{\alpha})$ অনুসৃত হয়; এবং
	\item $\alpha$ সীমা ক্রমসংখ্যা হলে $(\forall \beta \in
	\alpha)\phi(\beta)$ থেকে $\phi(\alpha)$ অনুসৃত হয়।
\end{enumerate}
তাহলে $\forall \alpha \phi(\alpha)$।
\end{thm}

\begin{proof}
প্রতিবিপরীত রূপটি প্রমাণ করি। তাই ধরি, কোনো ক্রমসংখ্যার
জন্য $\lnot\phi$ সত্য; তাদের মধ্যে ক্ষুদ্রতমটি
$\gamma$। তখন হয় $\gamma = \emptyset$; নয়
কোনো $\alpha$-এর জন্য $\gamma = \ordsucc{\alpha}$,
যেখানে $\phi(\alpha)$ সত্য; নয়তো $\gamma$ একটি
সীমা ক্রমসংখ্যা এবং $(\forall
\beta \in \gamma)\phi(\beta)$। প্রতিটি ক্ষেত্রেই
অনুমানগুলির বিরোধ হয়।
\end{proof}
\noindent
আরেকটি সংকেতলিপি পরে কাজে লাগবে:

\begin{defn}\ollabel{defsupstrict}
$X$ ক্রমসংখ্যাগুলির একটি সেট হলে $\supstrict(X) = \bigcup_{\alpha
\in X} \ordsucc{\alpha}$।
\end{defn}
\noindent
এখানে ``lsub'' বলতে ``ক্ষুদ্রতম কঠোর ঊর্ধ্বসীমা'' বোঝায়।\footnote{কিছু
বইয়ে এর জন্য ``$\text{sup}(X)$'' লেখা হয়। কিন্তু অন্য
বইয়ে ``$\text{sup}(X)$'' লেখা হয় ক্ষুদ্রতম
\emph{অকঠোর} ঊর্ধ্বসীমার জন্য, যা কেবল $\bigcup X$।
$X$-এর বৃহত্তম সদস্য $\alpha$ থাকলে এই দুটি ধারণা
আলাদা: ক্ষুদ্রতম \emph{কঠোর} ঊর্ধ্বসীমা
$\ordsucc{\alpha}$, আর ক্ষুদ্রতম \emph{অকঠোর}
ঊর্ধ্বসীমা শুধু $\alpha$।} নিচের ফলটি এই
সংকেতলিপির তাৎপর্য ব্যাখ্যা করে:

\begin{prop}
$X$ ক্রমসংখ্যাগুলির একটি সেট হলে $\supstrict(X)$
হলো $X$-এর প্রত্যেক ক্রমসংখ্যার চেয়ে বড় এমন
ক্ষুদ্রতম ক্রমসংখ্যা।
\end{prop}

\begin{proof}
ধরি, $Y = \Setabs{\ordsucc{\alpha}}{\alpha \in X}$;
তাহলে $\supstrict(X) = \bigcup Y$। ক্রমসংখ্যাগুলি
পরিযায়ী এবং ক্রমসংখ্যার প্রত্যেক সদস্যও ক্রমসংখ্যা।
তাই $\supstrict(X)$ ক্রমসংখ্যাগুলির একটি পরিযায়ী
সেট; \olref[basic]{corordtransitiveord} অনুযায়ী
এটি ক্রমসংখ্যা।

যদি $\alpha \in X$ হয়, তবে $\ordsucc{\alpha} \in Y$।
তাই $\ordsucc{\alpha}
\subseteq \bigcup Y = \supstrict(X)$, এবং সুতরাং
$\alpha \in
\supstrict(X)$। অতএব $\supstrict(X)$, $X$-এর
প্রত্যেক ক্রমসংখ্যার চেয়ে কঠোরভাবে বড়।

অন্যদিকে, $\alpha \in \supstrict(X)$ হলে
কোনো $\beta \in X$-এর জন্য $\alpha \in
\ordsucc{\beta} \in Y$; তাই $\alpha \leq
\beta \in X$। অতএব যেকোনো কঠোর
ঊর্ধ্বসীমায় সেই সদস্যও থাকবে। যেহেতু
সদস্যটি ছিল বিবেচ্য সেটের নির্বিচার সদস্য,
$\supstrict(X)$-ই $X$-এর \emph{ক্ষুদ্রতম} কঠোর
ঊর্ধ্বসীমা।
\end{proof}

\end{document}
